\documentclass{article}
\usepackage{amssymb}
\usepackage{amsthm}
\usepackage{amsmath,amssymb}
\usepackage{graphicx} % Required for inserting images

\title{6104_HW6}
\author{Odysseus Zhang}
\date{March 2026}

\begin{document}
\subsection*{1(a) Convexity}

We want to show that the log-likelihood function
\[
f(\theta)=\sum_{i=1}^n \left[y_i \log h(\theta^\top x_i) + (1-y_i)\log\bigl(1-h(\theta^\top x_i)\bigr)\right]
\]
is concave, where
\[
h(t)=\frac{1}{1+e^{-t}}.
\]

First, note that
\[
h(\theta^\top x_i)=\frac{1}{1+e^{-\theta^\top x_i}}.
\]
Hence
\[
f(\theta)
=
\sum_{i=1}^n
\left[
y_i \log\left(\frac{1}{1+e^{-\theta^\top x_i}}\right)
+
(1-y_i)\log\left(1-\frac{1}{1+e^{-\theta^\top x_i}}\right)
\right].
\]

We compute the gradient:
\[
\nabla f(\theta)
=
\sum_{i=1}^n
\left[
y_i \cdot \frac{1}{h(\theta^\top x_i)} \nabla h(\theta^\top x_i)
+
(1-y_i)\cdot \frac{1}{1-h(\theta^\top x_i)} \nabla\bigl(1-h(\theta^\top x_i)\bigr)
\right].
\]
Since
\[
\nabla h(\theta^\top x_i)=h(\theta^\top x_i)\bigl(1-h(\theta^\top x_i)\bigr)x_i,
\]
we obtain
\[
\nabla f(\theta)
=
\sum_{i=1}^n
\left[
y_i(1-h(\theta^\top x_i))x_i-(1-y_i)h(\theta^\top x_i)x_i
\right].
\]
Therefore,
\[
\nabla f(\theta)
=
\sum_{i=1}^n x_i\bigl(y_i-h(\theta^\top x_i)\bigr).
\]

Now we differentiate again. Since \(y_i\) is constant with respect to \(\theta\),
\[
\nabla^2 f(\theta)
=
-\sum_{i=1}^n x_i \bigl(\nabla h(\theta^\top x_i)\bigr)^\top.
\]
Using
\[
\nabla h(\theta^\top x_i)
=
h(\theta^\top x_i)\bigl(1-h(\theta^\top x_i)\bigr)x_i,
\]
it follows that
\[
\nabla^2 f(\theta)
=
-\sum_{i=1}^n h(\theta^\top x_i)\bigl(1-h(\theta^\top x_i)\bigr)x_i x_i^\top.
\]
Equivalently, since
\[
h(\theta^\top x_i)\bigl(1-h(\theta^\top x_i)\bigr)
=
\frac{e^{-\theta^\top x_i}}{(1+e^{-\theta^\top x_i})^2},
\]
we can also write
\[
\nabla^2 f(\theta)
=
-\sum_{i=1}^n
\frac{e^{-\theta^\top x_i}}{(1+e^{-\theta^\top x_i})^2}\, x_i x_i^\top.
\]

Now for any \(v\in\mathbb{R}^d\),
\[
v^\top \nabla^2 f(\theta) v
=
-\sum_{i=1}^n
h(\theta^\top x_i)\bigl(1-h(\theta^\top x_i)\bigr)\, v^\top x_i x_i^\top v.
\]
Since
\[
v^\top x_i x_i^\top v = (x_i^\top v)^2 \ge 0,
\]
and
\[
h(\theta^\top x_i)\bigl(1-h(\theta^\top x_i)\bigr)\ge 0,
\]
we have
\[
v^\top \nabla^2 f(\theta) v
=
-\sum_{i=1}^n
h(\theta^\top x_i)\bigl(1-h(\theta^\top x_i)\bigr)(x_i^\top v)^2
\le 0.
\]
Therefore,
\[
\nabla^2 f(\theta)\preceq 0.
\]
Hence \(f(\theta)\) is concave. Equivalently, \(-f(\theta)\) is convex. \qed

\subsection*{2(a) Gradient Calculation}

\textbf{(i)} For a general link function \(h\), the log-likelihood is
\[
f(\theta)=\sum_{i=1}^n \left[y_i \log h(\langle x_i,\theta\rangle)+(1-y_i)\log\bigl(1-h(\langle x_i,\theta\rangle)\bigr)\right].
\]
Differentiating term by term gives
\[
\nabla f(\theta)
=
\sum_{i=1}^n
\left[
y_i \frac{h'(\langle x_i,\theta\rangle)}{h(\langle x_i,\theta\rangle)}x_i
-
(1-y_i)\frac{h'(\langle x_i,\theta\rangle)}{1-h(\langle x_i,\theta\rangle)}x_i
\right].
\]
Factoring out \(x_i h'(\langle x_i,\theta\rangle)\), we get:
\[
\nabla f(\theta)
=
\sum_{i=1}^n
x_i\,h'(\langle x_i,\theta\rangle)
\left[
\frac{y_i}{h(\langle x_i,\theta\rangle)}
-
\frac{1-y_i}{1-h(\langle x_i,\theta\rangle)}
\right].
\]
Combining the two fractions,
\[
\frac{y_i}{h(\langle x_i,\theta\rangle)}
-
\frac{1-y_i}{1-h(\langle x_i,\theta\rangle)}
=
\frac{y_i-h(\langle x_i,\theta\rangle)}
{h(\langle x_i,\theta\rangle)(1-h(\langle x_i,\theta\rangle))}.
\]
Therefore,
\[
\nabla f(\theta)
=
\sum_{i=1}^n
x_i\,
\frac{h'(\langle x_i,\theta\rangle)}
{h(\langle x_i,\theta\rangle)(1-h(\langle x_i,\theta\rangle))}
\Bigl(y_i-h(\langle x_i,\theta\rangle)\Bigr).
\]

\textbf{(ii)} For the canonical logit link
\[
h(\gamma)=\frac{1}{1+e^{-\gamma}},
\]
we differentiate to get
\[
h'(\gamma)=\frac{e^{-\gamma}}{(1+e^{-\gamma})^2}.
\]
Also,
\[
1-h(\gamma)=1-\frac{1}{1+e^{-\gamma}}
=\frac{e^{-\gamma}}{1+e^{-\gamma}}.
\]
Hence
\[
h(\gamma)(1-h(\gamma))
=
\frac{1}{1+e^{-\gamma}}
\cdot
\frac{e^{-\gamma}}{1+e^{-\gamma}}
=
\frac{e^{-\gamma}}{(1+e^{-\gamma})^2}
=
h'(\gamma).
\]

\textbf{(iii)} Substituting \(h'(\gamma)=h(\gamma)(1-h(\gamma))\) into part (i), we obtain
\[
\nabla f(\theta)
=
\sum_{i=1}^n
x_i \bigl(y_i-h(\langle x_i,\theta\rangle)\bigr).
\]

\subsection*{2(b) Hessian Calculation}

From part (a)(iii), we have
\[
\nabla f(\theta)=\sum_{i=1}^n x_i\bigl(y_i-h(\langle x_i,\theta\rangle)\bigr).
\]
Hence the \(\ell\)-th component of the gradient is
\[
\frac{\partial f}{\partial \theta_\ell}
=
\sum_{i=1}^n x_i(\ell)\bigl(y_i-h(\langle x_i,\theta\rangle)\bigr).
\]
Now differentiate with respect to \(\theta_k\):
\[
\frac{\partial^2 f}{\partial\theta_k\partial\theta_\ell}
=
\sum_{i=1}^n
\frac{\partial}{\partial\theta_k}
\left[
x_i(\ell)\bigl(y_i-h(\langle x_i,\theta\rangle)\bigr)
\right].
\]
Since \(x_i(\ell)\) and \(y_i\) do not depend on \(\theta\), this becomes
\[
\frac{\partial^2 f}{\partial\theta_k\partial\theta_\ell}
=
-\sum_{i=1}^n x_i(\ell)\,
\frac{\partial}{\partial\theta_k}h(\langle x_i,\theta\rangle).
\]
By chain rule,
\[
\frac{\partial}{\partial\theta_k}h(\langle x_i,\theta\rangle)
=
h'(\langle x_i,\theta\rangle)
\frac{\partial}{\partial\theta_k}\langle x_i,\theta\rangle.
\]
Also,
\[
\langle x_i,\theta\rangle=\sum_{j=1}^d x_i(j)\theta_j,
\]
so
\[
\frac{\partial}{\partial\theta_k}\langle x_i,\theta\rangle=x_i(k).
\]
Therefore,
\[
\frac{\partial^2 f}{\partial\theta_k\partial\theta_\ell}
=
-\sum_{i=1}^n x_i(\ell)x_i(k)\,h'(\langle x_i,\theta\rangle).
\]
Since for the canonical logit link,
\[
h'(t)=h(t)(1-h(t)),
\]
we conclude
\[
\frac{\partial^2 f}{\partial\theta_k\partial\theta_\ell}
=
-\sum_{i=1}^n x_i(k)x_i(\ell)\,h(\langle x_i,\theta\rangle)\bigl(1-h(\langle x_i,\theta\rangle)\bigr).
\]

\subsection*{(c) Newton Update and Weighted Iterative Least Squares}
\subsection*{(i)}
Since we already have:
\[
\hat y_i=h(\langle x_i,\theta\rangle),
\qquad
\hat y=(\hat y_1,\dots,\hat y_n)^\top,
\]
and
\[
W=\operatorname{diag}\bigl(\hat y_1(1-\hat y_1),\dots,\hat y_n(1-\hat y_n)\bigr).
\]
Also, the design matrix is
\[
X=
\begin{bmatrix}
x_1^\top\\
x_2^\top\\
\vdots\\
x_n^\top
\end{bmatrix}.
\]

First, from part (a)(iii),
\[
\nabla f(\theta)=\sum_{i=1}^n x_i\bigl(y_i-h(\langle x_i,\theta\rangle)\bigr).
\]
Using \(\hat y_i=h(\langle x_i,\theta\rangle)\), we have
\[
\nabla f(\theta)=\sum_{i=1}^n x_i(y_i-\hat y_i).
\]
Now,
\[
\sum_{i=1}^n x_i y_i = X^\top y,
\qquad
\sum_{i=1}^n x_i \hat y_i = X^\top \hat y.
\]
Hence
\[
\nabla f(\theta)=X^\top y - X^\top \hat y = X^\top (y-\hat y).
\]

Next, from part (b),
\[
\frac{\partial^2 f}{\partial\theta_k\partial\theta_\ell}
=
-\sum_{i=1}^n x_i(k)x_i(\ell)\,\hat y_i(1-\hat y_i).
\]
On the other hand, the \((k,\ell)\)-entry of \(X^\top W X\) is
\[
(X^\top W X)_{k\ell}
=
\sum_{i=1}^n X_{ik}W_{ii}X_{i\ell}.
\]
Since
\[
X_{ik}=x_i(k),
\qquad
X_{i\ell}=x_i(\ell),
\qquad
W_{ii}=\hat y_i(1-\hat y_i),
\]
it follows that
\[
(X^\top W X)_{k\ell}
=
\sum_{i=1}^n x_i(k)\hat y_i(1-\hat y_i)x_i(\ell).
\]
Therefore,
\[
\frac{\partial^2 f}{\partial\theta_k\partial\theta_\ell}
=
-(X^\top W X)_{k\ell}.
\]
Since this holds for all \(k,\ell\), we conclude
\[
\nabla^2 f(\theta)=-X^\top W X.
\]
\textbf{(ii)} We want to show that the Newton direction
\[
d=-\bigl[\nabla^2 f(\theta)\bigr]^{-1}\nabla f(\theta)
\]
is the solution to the weighted least squares problem
\[
d=\arg\min_d \sum_{i=1}^n w_{ii}\bigl(z_i-\langle x_i,d\rangle\bigr)^2.
\]

From part (c)(i), we have
\[
\nabla^2 f(\theta)=-X^\top W X
\qquad\text{and}\qquad
\nabla f(\theta)=X^\top(y-\hat y).
\]
Therefore,
\[
d
=
-\bigl[\nabla^2 f(\theta)\bigr]^{-1}\nabla f(\theta)
=
-\bigl[-X^\top W X\bigr]^{-1}X^\top(y-\hat y).
\]
Since
\[
\bigl[-X^\top W X\bigr]^{-1}=-(X^\top W X)^{-1},
\]
it follows that
\[
d
=
(X^\top W X)^{-1}X^\top(y-\hat y).
\]
Also, by definition,
\[
z=W^{-1}(y-\hat y),
\]
so
\[
y-\hat y=Wz.
\]
Substituting this into the above expression gives
\[
d
=
(X^\top W X)^{-1}X^\top W z.
\]

Now consider the weighted least squares objective
\[
Q(d)=\sum_{i=1}^n w_{ii}\bigl(z_i-\langle x_i,d\rangle\bigr)^2.
\]
In matrix form, it can be written as
\[
Q(d)=(z-Xd)^\top W(z-Xd).
\]
We get
\[
Q(d)
=
z^\top W z-z^\top WXd-d^\top X^\top W z+d^\top X^\top W X d.
\]
Since \(z^\top WXd\) is a scalar, we have
\[
z^\top WXd=d^\top X^\top W z.
\]
Hence
\[
Q(d)
=
z^\top W z-2d^\top X^\top W z+d^\top X^\top W X d.
\]

Now differentiate with respect to \(d\):
\[
\nabla_d Q(d)
=
-2X^\top W z+2X^\top W X d.
\]
Setting the gradient equal to zero gives
\[
-2X^\top W z+2X^\top W X d=0,
\]
so
\[
X^\top W X d=X^\top W z.
\]
Assuming \(X^\top W X\) is invertible, we obtain
\[
d=(X^\top W X)^{-1}X^\top W z.
\]

This is exactly the same expression as the Newton direction. Therefore, the Newton direction is the solution to the weighted least squares problem
\[
d=\arg\min_d \sum_{i=1}^n w_{ii}\bigl(z_i-\langle x_i,d\rangle\bigr)^2.
\]


\subsection*{3. Computational Runtime}

\textbf{Evaluation of the log-likelihood \(f(\theta)\).}
\\ Since 
\[
f(\theta)
=
\sum_{i=1}^n
\left[
y_i \log h(\langle x_i,\theta\rangle)
+
(1-y_i)\log\bigl(1-h(\langle x_i,\theta\rangle)\bigr)
\right].
\]
For each \(i\), computing the inner product \(\langle x_i,\theta\rangle\) requires \(O(d)\) operations.  
Then evaluating \(h(\langle x_i,\theta\rangle)\), the logarithms, and the remaining scalar arithmetic requires \(O(1)\) operations.  
Thus each summand costs \(O(d)\), and since there are \(n\) summands, the total runtime is
\[
O(nd).
\]

\textbf{Evaluation of the gradient \(\nabla f(\theta)\).}
\\ Since
\[
\nabla f(\theta)=\sum_{i=1}^n x_i\bigl(y_i-h(\langle x_i,\theta\rangle)\bigr),
\]
for each \(i\), we first compute \(\langle x_i,\theta\rangle\), which costs \(O(d)\).  
Then computing \(h(\langle x_i,\theta\rangle)\) and the scalar \(y_i-h(\langle x_i,\theta\rangle)\) costs \(O(1)\).  
Finally, multiplying the vector \(x_i\in\mathbb{R}^d\) by this scalar and adding it to the running sum costs \(O(d)\).  
Hence each term costs \(O(d)\), and over all \(n\) observations the total runtime is
\[
O(nd).
\]

\textbf{Evaluation of the Hessian \(\nabla^2 f(\theta)\).}

Using
\[
\nabla^2 f(\theta)
=
-\sum_{i=1}^n
h(\langle x_i,\theta\rangle)\bigl(1-h(\langle x_i,\theta\rangle)\bigr)\,x_i x_i^\top,
\]
for each \(i\), computing \(\langle x_i,\theta\rangle\) costs \(O(d)\), and evaluating the scalar factor
\[
h(\langle x_i,\theta\rangle)\bigl(1-h(\langle x_i,\theta\rangle)\bigr)
\]
costs \(O(1)\).  
The outer product \(x_i x_i^\top\) is a \(d\times d\) matrix, and forming it requires \(O(d^2)\) operations.  
Scaling this matrix and adding it to the running Hessian also costs \(O(d^2)\).  
Therefore each summand costs \(O(d^2)\), and over all \(n\) observations the total runtime is
\[
O(nd^2).
\]

\textbf{Computation of the Newton direction.}

The Newton direction is
\[
d=-\bigl[\nabla^2 f(\theta)\bigr]^{-1}\nabla f(\theta).
\]
To compute it, we need:

\begin{itemize}
    \item the gradient, which costs \(O(nd)\),
    \item the Hessian, which costs \(O(nd^2)\),
    \item solving a linear system with a \(d\times d\) matrix, which costs \(O(d^3)\).
\end{itemize}

Hence the total runtime is
\[
O(nd)+O(nd^2)+O(d^3)=O(nd^2+d^3).
\]
\subsection*{4. Bonus Problems}

Recall that the negative log-likelihood is
\[
h(\theta)=-f(\theta).
\]
From the previous parts, we know that
\[
\nabla^2 f(\theta)=-X^\top W(\theta)X,
\]
where
\[
W(\theta)=\operatorname{diag}\bigl(\hat y_1(1-\hat y_1),\dots,\hat y_n(1-\hat y_n)\bigr),
\qquad
\hat y_i=h(\langle x_i,\theta\rangle).
\]
Hence
\[
\nabla^2 h(\theta)=-\nabla^2 f(\theta)=X^\top W(\theta)X.
\]

\textbf{(i) \(h\) is not \(m\)-strongly convex for any \(m>0\).}

Suppose, for contradiction, that \(h\) is \(m\)-strongly convex for some \(m>0\). Then for every \(\theta\),
\[
\nabla^2 h(\theta)\succeq mI_d.
\]
Equivalently, for every \(v\in\mathbb{R}^d\),
\[
v^\top \nabla^2 h(\theta)v \ge m\|v\|_2^2.
\]

Now fix any \(v\in\mathbb{R}^d\). Since
\[
\nabla^2 h(\theta)=X^\top W(\theta)X,
\]
we have
\[
v^\top \nabla^2 h(\theta)v
=
v^\top X^\top W(\theta)Xv
=
(Xv)^\top W(\theta)(Xv)
=
\sum_{i=1}^n \hat y_i(1-\hat y_i)\,(x_i^\top v)^2.
\]
For the logistic link,
\[
\hat y_i=\frac{1}{1+e^{-\langle x_i,\theta\rangle}},
\qquad
\hat y_i(1-\hat y_i)\to 0
\quad\text{as}\quad
|\langle x_i,\theta\rangle|\to\infty.
\]

Choose any nonzero vector \(u\in\mathbb{R}^d\), and let \(\theta=t u\) with \(t\to\infty\). Then for each \(i\),
\[
0\le \hat y_i(1-\hat y_i)\le \frac14,
\qquad
\hat y_i(1-\hat y_i)\to 0.
\]
Therefore,
\[
v^\top \nabla^2 h(tu)v
=
\sum_{i=1}^n \hat y_i(1-\hat y_i)\,(x_i^\top v)^2
\longrightarrow 0
\qquad\text{as } t\to\infty.
\]
But if \(h\) were \(m\)-strongly convex, we would have
\[
v^\top \nabla^2 h(tu)v \ge m\|v\|_2^2 >0
\]
for every \(t\), which is impossible since the left-hand side tends to \(0\). This contradiction shows that \(h\) is not \(m\)-strongly convex for any \(m>0\).

\bigskip

\textbf{(ii) Gradient Lipschitz property.}

We now prove that \(\nabla h\) is Lipschitz with
\[
L\le \frac14 \lambda_{\max}(X^\top X).
\]

For the logistic link, for every \(i\),
\[
0\le \hat y_i(1-\hat y_i)\le \frac14.
\]
Hence
\[
0\preceq W(\theta)\preceq \frac14 I_n.
\]
Multiplying by \(X^\top\) on the left and \(X\) on the right gives
\[
0\preceq \nabla^2 h(\theta)=X^\top W(\theta)X \preceq \frac14 X^\top X.
\]
Therefore,
\[
\|\nabla^2 h(\theta)\|_2
\le
\left\|\frac14 X^\top X\right\|_2
=
\frac14 \|X^\top X\|_2
=
\frac14 \lambda_{\max}(X^\top X).
\]

Now let \(\theta,\theta'\in\mathbb{R}^d\). By the mean value theorem for vector-valued functions,
\[
\nabla h(\theta)-\nabla h(\theta')
=
\left(\int_0^1 \nabla^2 h\bigl(\theta'+t(\theta-\theta')\bigr)\,dt\right)(\theta-\theta').
\]
Taking norms and using the bound above,
\[
\|\nabla h(\theta)-\nabla h(\theta')\|_2
\le
\int_0^1
\left\|\nabla^2 h\bigl(\theta'+t(\theta-\theta')\bigr)\right\|_2\,dt \;\|\theta-\theta'\|_2
\]
\[
\le
\int_0^1 \frac14 \lambda_{\max}(X^\top X)\,dt \;\|\theta-\theta'\|_2
=
\frac14 \lambda_{\max}(X^\top X)\,\|\theta-\theta'\|_2.
\]
Thus \(\nabla h\) is Lipschitz, with Lipschitz constant
\[
L\le \frac14 \lambda_{\max}(X^\top X).
\]

Therefore, \(h\) is not \(m\)-strongly convex for any \(m>0\), but it is \(L\)-gradient Lipschitz with
\[
L\le \frac14 \lambda_{\max}(X^\top X).
\]
\end{document}
